\documentclass[11pt,letterpaper]{article}
\usepackage[T1]{fontenc}
\usepackage[utf8]{inputenc}
\usepackage{lmodern,microtype}
\usepackage[margin=0.9in]{geometry}
\usepackage{amsmath,amssymb,amsthm,mathtools}
\usepackage{booktabs,longtable,array,tabularx}
\usepackage{xcolor,xurl,hyperref,enumitem,fancyhdr}
\hypersetup{colorlinks=true,linkcolor=blue!45!black,urlcolor=blue!45!black,
pdftitle={PA194P5: a five-register permutation and a comparison of PA consistency machines},
pdfauthor={Research prepared for Vladimir Reshetnikov}}
\urlstyle{tt}
\setlength{\emergencystretch}{3em}
\setlength{\parindent}{0pt}
\setlength{\parskip}{5pt}
\setlength{\headheight}{14pt}
\setlist{nosep,topsep=5pt}
\pagestyle{fancy}\fancyhf{}
\fancyhead[L]{\small PA194P5 and the 194-state frontier}
\fancyhead[R]{\small 14 November 2026}
\fancyfoot[C]{\thepage}
\newcommand{\PA}{\mathrm{PA}}
\newcommand{\Con}{\operatorname{Con}}
\newcommand{\Halt}{\operatorname{Halt}}
\newcommand{\Prf}{\operatorname{Prf}}
\newcommand{\file}[1]{\texttt{\nolinkurl{#1}}}
\newcommand{\old}{M_{194}}
\newcommand{\new}{M_{\mathrm{P5}}}
\newcommand{\snapshot}{\texttt{4be5c6e3a10292a7466f5d3687e3b39615ff049b}}
\newtheorem{theorem}{Theorem}[section]
\newtheorem{lemma}[theorem]{Lemma}
\newtheorem{proposition}[theorem]{Proposition}
\title{\textbf{PA194P5}\\A five-register permutation and a comparison\\of PA consistency machines}
\author{Research prepared for Vladimir Reshetnikov}
\date{14 November 2026}
\begin{document}
\maketitle
\begin{abstract}
This report supplies a deterministic binary, one-tape machine with
194 working states, obtained by relocating five logical registers in the
primary PA194 implementation. Seven transition targets change; all other
381 transitions retain their original dense numbering and actions. The
logical source program, numerical calculus, and decoder grammar are
unchanged. The new table is distinct from every 194-control transition
graph in the inspected current corpus, even after allowing state
renaming, a different graph root, tape reflection, symbol interchange,
and erasure of the actions on transitions into halt.

The physical simulation is checked for all 4,780 symbolic source cases,
with arbitrary register values and the symbolic header gap. A complete
blank-tape execution reaches precisely the parent's clean source frame
after 56,158 transitions, compared with 58,154 for the parent. The
permutation therefore gives an ordinary PA-internal finite-trace
equivalence with the supplied primary PA194 machine, and inherits that
packet's written PA consistency calibration. This is a transport of an
existing calibration, not a new end-to-end Lean or Rocq calibration.

The requested strict reduction below 194 states was not obtained.
The delivered result is the permitted same-size physical-design
fallback. The comparison pins all sources to one revision, covers
61 placed research packets and all eleven current incoming archives,
and separates numerical table counts from proof status. We also give
a concrete obstruction to repairing a supplied 193-state runtime core
by physical register permutations alone.
\end{abstract}
\clearpage
\tableofcontents
\clearpage

\section{Result, scope, and proof boundary}

The final table is \file{machine/pa194p5.tm}. It has working controls
$0,\ldots,193$, start control $0$, and terminal halt label $194$.
Every working control has exactly two transitions, one for each read
symbol. The tape alphabet is $\{0,1\}$; the tape is two-way infinite and
initially all zero; the head starts at coordinate zero. Every transition
moves one cell left or right. The halt label has no transition rows.
There are 388 transitions, and all 194 working controls are reachable
in the pointed transition graph.

The main theorem proved by the new construction is
\[
             \PA\vdash
             \bigl(\Halt(\new)\leftrightarrow\Halt(\old)\bigr).
\]
Composing this transport with the supplied primary packet's written
calibration gives
\[
             \PA\vdash
             \bigl(\neg\Halt(\new)\leftrightarrow\Con(\PA)\bigr).
\]
The second assertion uses the inherited numerical-calculus calibration.
The new computations and proof below establish the physical transport;
they do not independently replace all logical lemmas in that
calibration. Section~\ref{sec:calibration} records its content and
dependencies. Neither displayed assertion is presented as a newly
compiled proof-assistant endpoint.

This distinction matters because the repository has a separate,
stronger project: internalizing the consistency equivalence in iPRA.
Its README explicitly says the machine-specific iPRA directions remain
open. A standard finite-run Lean theorem, a Python certificate, and a
conditional iPRA assembly have different logical scopes
\cite{status}. The present report addresses the requested ordinary PA
provability statement by preservation of the supplied source program.

The snapshot is ProveIt commit
\begin{center}\small\snapshot.\end{center}
All novelty and corpus claims refer to this revision, not to an
unbounded claim about the literature or later repository changes.
The comparison accounts for archived experiments, including rejected
tables, rather than treating every filename as an established result.

\section{What the comparison shows}

\subsection{Corpus accounting}

There are 61 placed packets under
\file{Computability/BusyBeaver/Research/PA/docs} and eleven ZIP archives
under \file{docs/incoming}. The placement ledger records 7,130
historical table occurrences in its 61 packets, 786 distinct file
contents and 535 classes under state renaming. Most duplicate table
files were removed at placement and replaced by pointers. These
historical counts therefore differ from a count of files present in
the current tree \cite{catalog,placement}.

The fresh census enumerates all current \file{*.tm} paths in the PA
tree, verifies their bytes against the Git blob identities, and
recursively opens all eleven incoming ZIPs. The package includes
the occurrence ledger, literal tables, graph representatives,
archive digests, exclusions and an offline replay of the frontier
comparison. The final machine-readable summary is
\file{comparison/census.json}.
\begin{center}
\begin{tabular}{@{}lr@{}}
\toprule Current-corpus quantity & Count\\\midrule
Direct PA-tree \file{.tm} paths &548\\
Incoming archives opened &11\\
Historical PA archive paths opened &6\\
Distinct recursive archive payloads &25\\
Literal table occurrences &3,334\\
Distinct literal byte strings &762\\
Spanning action-graph classes, all sizes &638\\
Spanning 194-state graph classes &82\\
Additional nonspanning 194-control diagnostics &2\\
193-state graph classes, all rejected or incomplete &53\\
Spanning graph classes with at least 195 controls &493\\
\bottomrule
\end{tabular}
\end{center}
The other spanning classes are tiny diagnostic machines. Six literal
identities have no spanning root: the two 194-control diagnostics and
four smaller control fragments. Twenty-two \file{.tm} occurrence paths
contain placeholders or partial symbolic listings rather than complete
numeric tables; their exact classifications are retained in
\file{comparison/unparsed_diagnostics.json}. The literal census parses numerical rows in \file{.tm/.txt/.tsv/.csv}
files; it does not enumerate every table expressible by arbitrary JSON
encodings or source generators. The excluded literal files contain no
unclassified complete numeric 194-state candidate.


The graph census deliberately ignores the declared start when comparing
unpointed action graphs: several retained native tables are numbered
densely even though state zero is not their intended start. This avoids
counting a renumbering with a mistaken start convention as a new
design. For a complete table with a spanning root, traversing from
each possible spanning root and reading the two branches in a fixed
order gives an exact finite canonical comparison. The stronger
novelty check also permits reflection and global symbol interchange,
and ignores the write and movement of transitions into halt.
PA194P5 matches none of the 82 current spanning 194-state representatives
under these identifications. Two further 194-control encoding diagnostics
have no spanning root (their largest reachable components have only three
controls); this invariant excludes an isomorphism with PA194P5. Their action
histograms also distinguish them from each other under all four global
symmetries. Thus all 84 recognized 194-control graph classes are accounted
for. The two diagnostics are not proposed blank-start consistency machines.

Of the 82 spanning frontier classes, the attributed inventory categories
are 40 principal constructions, seven supplementary constructions,
five source-validated research rotations, 25 research auxiliaries, and
five rejected tables. This is not a list of 82 independently certified
PA consistency machines.

These are graph and file comparisons. They do not assert that every
listed graph has a correct PA calibration, that every state occurs
on the blank run, or that differently shaped graphs have different
halting predicates. The 72-entry packet ledger in
Appendix~\ref{app:ledger} distinguishes reported contributions and
duplicate names. It is a complete packet inventory, not a fresh
proof audit of every historical manuscript.

\clearpage
\subsection{Architectural families}

\begin{longtable}{@{}p{0.21\linewidth}p{0.73\linewidth}@{}}
\toprule Family & Substantive difference\\\midrule\endhead
U208 families &
Conditional, entry-context, history, and three-cell reductions of
the older PA209 implementation; larger comparison baselines.\\[4pt]
PA196 and source experiments &
Imported register enumerator and binary interpreter, plus numerous
tree re-associations, source abbreviations and unsuccessful lower-state
candidates.\\[4pt]
PA195 and 195S &
Original diagonal and separate-head variants. Some use 68 readers
with seven folds; others use 67 readers with six folds.\\[4pt]
Primary PA194 &
The 190-token separate-head emitter, with 134 fixed controls and
$67-7$ reader controls after simultaneous folds.\\[4pt]
Alternate PA194 &
A different $68-8$ allocation with the same total of 194.\\[4pt]
C, CS, N, H, HS, SUMS &
Changes to successful modus ponens, dummy stack pushes, the cached
head, and the Shift contradiction test. The stack invariant is
part of their proof obligation.\\[4pt]
Q, D, and coefficient families &
Curried quantifier-premise exchange; a different separately emitted
diagonal instance; or coefficient-five/six formula coding, requiring
explicit formula and proof recoding.\\[4pt]
J, ZR, A, F, R &
Known-zero decrements used as jumps, changes to source entry or AX
return, and jointly re-associated decoders. Their source invariant
must also be reconciled with the blank startup.\\[4pt]
AX11 and recent combinations &
An eleven-entry axiom interface and combinations such as QH5,
QHS, QN5 and QN6. Entry counts change selector words; names alone
do not distinguish new graphs.\\[4pt]
PA194P5 (this report) &
A three-cycle of physical slots $2,3,4$ and exchange of slots $7,8$.
The theorem enumerator, source instructions and logical axioms are
retained; the represented tape store and physical execution change.\\
\bottomrule
\end{longtable}

Several names are aliases. The principal C graph appears in multiple
packets; QN and NQ name the same principal graph; one NS is the HS
graph, whereas another NS uses a symmetric \emph{sum} guard and is
different. CSZ and X are also aliases. The comparison files preserve
the provenance rather than flattening these names into a single
unqualified list of purported improvements.

\section{Exact definition of PA194P5}

Let the source registers be
\[
                    (N,T,W,A,D,H,X,Y,L).
\]
The parent assigns them to physical file positions
\[
                    p=(6,2,0,4,3,7,8,5,1).
\]
Define the permutation of physical positions
\[
       \pi=(0,1,3,4,2,5,6,8,7)
\]
in image-list notation: $\pi(2)=3$, $\pi(3)=4$, $\pi(4)=2$,
$\pi(7)=8$, and $\pi(8)=7$. The new assignment is
\[
                    p'=\pi\circ p=(6,3,0,2,4,8,7,5,1).
\]

\begin{center}
\begin{tabular}{@{}lrrrrrrrrr@{}}
\toprule Logical register & $N$&$T$&$W$&$A$&$D$&$H$&$X$&$Y$&$L$\\
\midrule Parent position &6&2&0&4&3&7&8&5&1\\
New position &6&3&0&2&4&8&7&5&1\\
\bottomrule
\end{tabular}
\end{center}

The five moved registers are $T,A,D,H,X$. This changes the order in
which unary data blocks occupy the tape. It is a physical store
permutation, with seven instruction-router changes. It is not a
new logical calculus or a new asymptotic proof-search algorithm.

In the \emph{parent's existing dense numbering}, replace precisely
the following seven rows. Every row in the table below reads one,
writes zero, and moves right.
\begin{center}
\begin{tabular}{@{}rrr@{}}
\toprule Dense control & Old target & New target\\\midrule
42&60&85\\
48&68&114\\
73&85&60\\
109&95&68\\
110&95&68\\
120&68&114\\
124&114&95\\
\bottomrule
\end{tabular}
\end{center}
All other 381 rows retain their parent labels and actions.
Appendix~\ref{app:table} gives the complete final table.

The parent dense file has SHA-256
\begin{quote}\small\nolinkurl{4bfea62b17aad7d8c8c5e682fe8942e0a86b93e4351f8e675e32705bdaae007f}.\end{quote}
The delivered final dense file has SHA-256
\begin{quote}\small\nolinkurl{cf3a2c2d3ba8e95d5c0ea825c3f279ae529b6591cc1280264b0ba3793e77d026}.\end{quote}
The native file, start 74 and halt 254, has SHA-256
\begin{quote}\small\nolinkurl{a9d3a13786b6a0018713d8803865860b5e9e26692fde34975545c3bb0f09501c}.\end{quote}
The fixed dense numbering is intentional: it makes the seven changes
directly reviewable. The included dense-to-native label list verifies
the complete conjugacy of the two presentations.

\section{Why the register relocation preserves the computation}
\label{sec:transport}

\subsection{Source frames and physical stores}

A natural register value $v$ is represented by a block $0\,1^{v+1}$.
A source frame consists of the left marker, the source-address bits,
a zero gap, the right delimiter, nine work blocks, and a nonempty
spare tail. At a source address $a$, the right delimiter is at
distance $48+n$ from the frame head, where $n$ is an arbitrary
natural number. Thus the zero gap has length $48+n-|a|$.

For a logical store $v=(v_0,\ldots,v_8)$, write
$\mathcal E_p(a,n,v,s)$ for the complete tape encoding with
$v_r$ at physical position $p(r)$ and spare tail $s$.
Define $\mathcal E_{p'}$ by the new assignment. The exterior zero
tape, source bits, delimiters and spare tail are retained. In
particular, the relation between encodings is a permutation of
whole unary blocks, not a pointwise relabeling of tape coordinates.

The physical selectors, in native control labels, are
\[
 (S_0,\ldots,S_8)=(51,36,25,48,32,33,22,55,60).
\]
Every instruction-router transition entering one of these selectors
from outside that list has the form $0R S_j$ on reading one.
Replacing its target by $S_{\pi(j)}$ gives precisely the seven
changes above; unchanged targets correspond to fixed positions
of $\pi$.

\begin{lemma}[Store operations commute with relocation]
For every logical register $r$, increment, zero test and positive
decrement on the representation $\mathcal E_{p'}$ realize exactly
the same logical operation as on $\mathcal E_p$. The remaining
logical registers are preserved, and the inherited spare-tail and
header effects agree.
\end{lemma}
\begin{proof}
The selector $S_j$ chooses physical block $j$. The new router for
register $r$ chooses $S_{\pi(p(r))}=S_{p'(r)}$, the block containing
$v_r$. The mark, increment and decrement routines themselves are
unchanged. Their contracts are uniform in the selected position,
the values of all work registers, and the finite nonempty spare
list. Increment increases the selected value; a zero test reports
whether that value is zero; positive decrement lowers it by one.
All other represented values remain fixed. The increment contract
includes its spare-list extension, so that effect is not silently
discarded. The address and gap updates are the same control
operations as in the parent. These identities prove the claim for
arbitrary values, not for a bounded sample of stores.
\end{proof}

The complete arithmetic-box state set used by the retained checker
has 23 controls. Every one of their 46 transition rows is bytewise
unchanged in PA194P5. The selector-target changes occur outside
this set. The checker reuses the four inherited register-box
contracts with exactly their old row premises. The actual provider of
these lemmas is the upstream theorem \file{M6.Native.boxLemmas}, not
merely the hypothesis structure in \file{Runtime.lean}. It requires
agreement on 44 engine rows and the append-one tail
$198/0=1L40$; preserving all 46 rows is stronger than these premises.
The generic proof sources and their frozen reference table are retained
under \file{reference/lean/} \cite{boxlemmas}. They are dependencies of
the mathematical argument, not a newly compiled PA194P5 Lean endpoint.

\begin{lemma}[Primitive source simulation]
The 3,673-instruction source program is the same in the parent
and PA194P5. From every valid frame $\mathcal E_{p'}$, each
nonhalting source instruction has a terminating, positive-length
physical simulation leading to the correct successor frame and
store. A source halt reaches the physical halt.
\end{lemma}
\begin{proof}
The decoder grammar, every primitive opcode and operand, and every
resolved source successor are unchanged. At a register operation,
apply the preceding lemma. All remaining address-decoder, return
and structured-jump transitions are retained from the parent.
The finite verification has 4,780 cases: one per increment or
unconditional jump, both decrement outcomes, and the halting case.
Each symbolic frame retains arbitrary values and the affine gap
parameter. The completed symbolic traces establish the intended
successor address, whole represented store and complete frame;
a positive physical transition is required before a contract
can close.
\end{proof}

The symbolic verifier is an executable certificate checker with
ordinary mathematical loop and register lemmas. It is not a
proof assistant. The proof-producing interpretation is elementary:
a constant-symbol scan is proved by induction on its length,
and a register-file operation by induction on the number and
length of the blocks it traverses. Every certificate is then a
finite composition of those lemmas and literal transitions.
Such finite-trace constructions are primitive recursive and
their correctness is provable in PA.

\subsection{Blank initialization is a separate obligation}

Changing a register allocation can preserve every well-formed
source contract while destroying blank initialization. We tested
this obligation independently, without assuming that the source
invariant holds before a clean source frame exists.

The parent reaches source site 25, native control 84 (dense 24),
head $-2$, with the logical vector $(2,0,\ldots,0)$ after
58,154 steps. PA194P5 reaches the \emph{identical complete
configuration} after 56,158 steps. The source path is $100000$;
the header distance is 48; the nine physical values are
$(0,0,0,0,0,0,2,0,0)$; there are 26 one-valued spare blocks.
Both tapes have exactly the same 66 one-cells and are zero
elsewhere. The complete support is recorded in the replay receipt.

This equality is consistent with the new representation because
$\pi$ fixes position six, which contains the only nonzero work
register in the seed. The 1,996-step saving concerns initialization
only; no general speedup of the subsequent proof search is claimed.

\begin{theorem}[PA-internal finite halting transport]
\[
             \PA\vdash
             \Halt(\new)\leftrightarrow\Halt(\old).
\]
\end{theorem}
\begin{proof}
Let $R$ be their common deterministic register program from site 25
and the clean logical seed. Given a finite halting source trace,
prepend the appropriate verified blank initialization and concatenate
the positive finite simulation segments of its instructions.
This constructs a finite halting tape trace of either machine.

Conversely, suppose one of the physical runs first halts after $T$
steps. Iterate the total source step function for $T+1$ instructions,
with source halt treated as an absorbing event. If no source halt
has occurred, the positive physical segments for those instructions,
together with initialization, give a nonhalting physical execution
extending beyond time $T$. Determinism contradicts the alleged
first physical halt. Thus that bounded source iteration contains
a halt. This supplies both physical-to-source reductions.

All statements concern coded finite traces. Source iteration,
register-block scanning and the finite concatenations are primitive
recursive; induction in PA proves their correctness. No assumption
of actual PA consistency, and no reflection principle for PA,
is used in this transport.
\end{proof}

\section{Inherited consistency calibration}
\label{sec:calibration}

For clarity, we spell out the logical dependency which the physical
transport preserves. This section summarizes the supplied
separate-diagonal calibration \cite{primary,calibration}; its retained
source is \file{reference/primary_calibration.tex}.
It is the same calibration relied on by the primary PA194 packet,
not a newly mechanized logical proof in this report.

For Cantor pairing $p(a,b)=(a+b)(a+b+1)/2+b$, its numerical
formula constructors are
\[
 I(A,B)=4p(A,B),\quad M(a,b)=4p(a,b)+1,\quad
 U(a,B)=4p(B,a)+2,\quad E(a,b)=4p(a,b)+3.
\]
Code zero is truth and code one is the contradiction $M(0,0)$.
Negation is $N(B)=I(B,1)$. The transposition in the quantifier
constructor is essential.

The separate diagonal clause is
$\Gamma(B)=B\to U(B,B)$. Every variable label of a nonzero
formula code is strictly below that code; truth is handled
separately. Thus the variable labeled by $B$ is absent from its
own formula and the diagonal clause is a valid vacuity instance.
The AX command also supplies the guarded gate
\[
\begin{aligned}
 V(a,b,c)&=M(a,b)\to U(c,M(a,b)),\\
 P(a,b,c)&=c\to M(c,b),\\
 S(a,b,c)&=U\bigl(b,(M(a,b)\to P(a,b,c))
                       \wedge(M(c,b)\to c)\bigr),\\
 Q(a,b,c)&=U(c,S(a,b,c)\to N(P(a,b,c))),\\
 G(a,b,c)&=N(V(a,b,c)\to Q(a,b,c)),
\end{aligned}
\]
with $c\ne a,b$, and truth at the guarded position otherwise.
The other six axiom heads, the dominated implicational head,
modus ponens, generalization and the ten-entry stack protocol
are retained without alteration.

For the source-to-PA direction, interpret membership by finite
binary membership $\operatorname{Bit}(x_a,x_b)$, equality by
arithmetic equality, and $U$ by universal quantification. The bit
convention is
\[
 \operatorname{Bit}(s,z)\;\longleftrightarrow\;
 \exists u,v\bigl(s=(2u+1)2^z+v\;\wedge\;v<2^z\bigr).
\]
PA proves $\operatorname{Bit}(s,z)\to z<s$, and hence refutes
the translation of code one. For any arithmetic formula
$\theta(z,\vec u)$, PA proves finite binary comprehension:
\[
 \forall n\,\exists s<2^n\,\forall z<n\,
   \bigl(\operatorname{Bit}(s,z)\leftrightarrow\theta(z,\vec u)\bigr).
\]
A proof is generated by the usual induction, retaining the old
binary code or adding $2^n$ according as $\theta(n,\vec u)$
is false or true. This is a uniform syntactic proof construction,
not a uniform truth predicate for arithmetic. It supplies the
pointed-separation instances, with the coincident-label case
handled separately.

The source store is annotated by PA proofs of the translations
of both visible theorem codes and every live stack entry. AX,
MP and Shift preserve these annotations. Canonical zero-stack
collisions and the original $t+h=1$ halt test require both visible
proofs. A finite source halt yields a finite PA refutation.

For the PA-to-source direction, the inherited compiler recovers
the necessary two-way membership rewriting, neighborhood equality
and equality-witness interfaces from the diagonal and guarded
separation schemas. It then translates a PA refutation into a
finite command derivation. Exhaustive command scheduling finds
that derivation. The stack selectors must use the ten-entry
separate-head protocol; the nine-entry predecessors cannot be
substituted unchanged. In particular, the supplied proof gives
the adjusted pop, push, discard and generalization words.

The resulting inherited syntactic calibration is
\[
 \PA\vdash
 \Halt(R)\leftrightarrow
 \exists e\,\Prf_{\PA}(e,\ulcorner0=1\urcorner).
\]
Combining it with the physical transport yields the stated
consistency equivalence for PA194P5. The unchanged source is
the reason the new construction incurs no new axiom-soundness,
completeness, proof-enumeration or stack-selection obligation.
A reader requiring a fully kernel-checked calibration must
still address the open endpoints documented in the repository.

\section{Reproducible verification}

The Python replay \file{verify.py} uses only the standard library.
It binds the parent by SHA-256, reconstructs the seven changes,
derives the native table via the retained label conjugacy, checks
the physical permutation against all selector-entry edges,
checks graph reachability and the unique halt, verifies every
arithmetic-box row, and replays the complete blank startup and
all source contracts.

\begin{center}
\begin{tabular}{@{}lrr@{}}
\toprule Quantity & Primary PA194 & PA194P5\\\midrule
Working controls &194&194\\
Transitions &388&388\\
Source instruction sites &3,673&3,673\\
Symbolic source cases &4,780&4,780\\
Symbolic steps, total &1,479,348&1,479,348\\
Largest symbolic case &497&497\\
Blank startup steps &58,154&56,158\\
One-cells in clean seed &66&66\\
\bottomrule
\end{tabular}
\end{center}

The unchanged symbolic step count is expected: the register-box
operations are contracted by their uniform lemmas. It is not a
claim that the physical running times on all registers are the
same. The literal blank interpreter counts actual tape steps.

Independent JavaScript and Python executions agree on the complete
source-contract result. A second independently written Set-based
interpreter reproduced the whole blank seed. Graph comparison
uses the literal tables, independently of the source compiler.
This separation of representations checks meaningful risks:
an erroneous target edit, a dense/native confusion, an incorrect
register interpretation, an omitted spare-tail effect, or a
bootstrap which never reaches a valid source frame.

\section{The search below 194 and why the apparent shortcuts fail}

\subsection{A precise negative local search}

For both primary PA194 and PA194P5 we collected symbolic contexts
from all source cases and the complete startup. The 171 non-box
controls have 29,347 and 29,325 distinct radius-six patterns,
respectively. Each pattern is a 13-cell tape window; arbitrary
register-box contents are expanded conservatively. The collection
observed 593,217 distinct symbolic configurations for each machine.

A local replacement could use up to 120 transitions while preserving
the whole local tape and leaving its exterior untouched. The
implemented searches found no contextual pair fold or deletion
using a head offset of two cells. On PA194P5, independent incoming-edge
bypasses failed for all 164 eligible controls. There are 334
individually equivalent single-row rewrites in the tested family;
none enabled a direct alias deletion or a subsequently synthesized
contextual pair fold. Arithmetic-box controls were not mutation
targets in this search.

The code, data and exact receipts are in
\file{research/local_reductions/}. These finite negatives are not
a lower bound on the number of states of a PA consistency machine.

\subsection{Accounting for all 53 supplied 193-control graphs}

The complete census finds 53 distinct 193-control graphs. The following
statuses are recorded in their source packets, with exact evidence paths
in \file{comparison/rejected193.json}; the table does not claim a fresh
replay of every historical search.
\begin{center}
\begin{tabular}{@{}lrp{0.50\linewidth}@{}}
\toprule Family & Graphs & Recorded obstacle\\\midrule
Coefficient three &1& Logical refutation of the searched calculus\\
Original runtime core &1& Source contracts pass; blank startup fails\\
Append-zero core &1& Failed initialization experiment\\
AX11 composition &1& All 193 starts halt or enter outward rays\\
QN6 composition &1& Native start 74 enters an outward ray\\
ZFC-style relay &1& Nine source-contract failures; blank outward ray\\
Other marker quotients &47& All-start screens halt or enter outward rays\\
\bottomrule
\end{tabular}
\end{center}
The marker-transplant archive has 49 tables, two of which repeat the
original and append-zero cores. Its 9,457 recorded start screens give
100 halts and 9,357 certified rays, with no accepted source frame and
no run left merely bounded. None of these tables carries a successful
blank-start PA consistency calibration. The next two subsections explain
the logical and physical obstructions separately.

\subsection{Rejected coefficient-three encoding}

The placed NQ packet contains an explicit 193-state table whose
receipt says that the coefficient-three calculus derives a
contradiction \cite{rejected193}. Its 4,719 symbolic physical
contracts pass. The physical compiler is therefore not the
reason to reject it: the source theory is unsound for the desired
calibration. Counting states and testing simulations alone would
mistakenly certify this example.

\subsection{The 193-state runtime core and its startup obstruction}
\label{sec:core193}

A different incoming experiment has 193 controls and passes all
4,780 source contracts, but is not a blank-start consistency
machine \cite{runtime193}. In native labels it identifies 301
with 88, redirects targets accordingly, and changes
\[
            (88,0)\mapsto0L51,\qquad
            (51,0)\mapsto0R122.
\]
At a well-formed unmarked register delimiter, the next bit is one.
The old $16/1$ and new $122/1$ actions are both $0L1$, which
transports the modified arithmetic marking interface in that
context. Blank initialization does not satisfy this next-bit-one
premise.

There is a short obstruction to fixing this core merely by
permuting the physical work registers. Before the first
selector-router call, the blank prefix is independent of that
permutation. The first relevant call is $45/1$ at time 78.
It erases the last rightmost one, at cell four, and enters
some selector $S_j$ to its right. All remaining cells to the
right are zero. Every possible $S_j$ runs through the zero
selector chain to $51/0$, then the new helper sends control
through $122/0$, $137/0$, $30/0$, $45/0$ and $24/0$ to 35.
Control 35 writes ones and moves right forever on this blank
ray. There is no later source initialization or halt.

Thus all nine possible selector destinations at that first call
fail, independently of later router choices. This rules out
every physical-register permutation as a repair with the same
bootstrap. The incoming QN6 packet had already computationally classified all
$9! = 362{,}880$ selector permutations as outward rays. The argument
here gives a short analytical certificate for that negative result;
it is not a claim of new empirical search coverage. The packet contains the finite-prefix
certificate in \file{research/architecture/} and the precise scope of this obstruction; it says
nothing against changing the startup mechanism or the helper
design itself.

The outward-ray induction is itself provable in PA. Consequently,
assuming PA is consistent, these exact blank-start variants cannot
also admit a PA proof that their nonhalting is equivalent to
$\Con(\PA)$: that would give PA a proof of its own consistency.
This uses the second incompleteness theorem. A mere standard-model
observation of nonhalting would not by itself make this distinction.

Two further bounded one-row searches used all 772 existing-control
write/move/target choices. A marking-interface search retained only
the original helper and the published reduced helper. An increment
search, allowing an arbitrary valid nonempty spare tail, retained only
the original $74/0=1R0$ row. These searches have explicit step and
spatial bounds in the architecture packet; unlike the ray proof,
they do not establish an unrestricted impossibility theorem.

\subsection{Two logical and decoder avenues}

The point-separation component alone tolerates a collision $c=a$
in a suitable empty-or-singleton argument. The actual gate also
asserts membership vacuity $V$, so simply deleting that guard is
invalid. With $c=a=0$, $b=1$, $x_0=1$ and $x_1=0$,
$\operatorname{Bit}(1,0)$ is true while
$\forall x_0\,\operatorname{Bit}(x_0,0)$ is false. This gives
a concrete counterexample to the tempting weakened gate.

Relocating the diagonal instance to $\Gamma(H)$ is a valid
schema-family idea, already present as an unsuccessful source
seed in an older packet. Two retained best grammars score
$134+70-7=197$, not 194. We also examined reordering independently
emitted heads and moving the diagonal block, with 21,151
bounded reassociation iterations. No competitive candidate
was found in those experiments. Reordering heads would require
new stack selectors and a new resolved-source check even when
the set of axioms is unchanged.

\section{Relation to the ZFC and concise-iPRA suggestions}

The ZFC tree gives concrete precedents for three routes examined here:
source reassociation, contextual state folds, and register-interface
specialization. Its K215 construction combines a command-local field reader,
command ordering, and a partial selector; its reported reallocation removes
one control while changing 35 retained rows. Conversely, the K/S-calculus
investigations explicitly distinguish a smaller source DAG from a smaller
complete binary machine \cite{zfc}. The relevant cost is the allocated
table including initialization and helper routines. The 193-state runtime
experiment above is an especially direct example of why the bootstrap cost
cannot be omitted.

The concise-iPRA project measures complete standalone proof checkers by
tokens, AST nodes and normalized source bytes, rather than Turing-machine
states. Its native scalar core derives logical structure from equation,
$Q$ and $D$ formulas. The supplied developments combine evaluation with
discharge, replace two introductions by one, and reduce compound transport
to a one-term equality rule, with explicit equivalence theorems
\cite{concise}. These are useful examples of reducing a checker interface
by deriving rules instead of merely deleting them. They do not provide a
priced binary compilation of a PA-refutation search.

Replacing the present searched calculus by iPRA would also change the
consistency claim unless accompanied by an explicit PA-provable translation
of refutations in both directions. The guarded-vacuity counterexample above
shows, within the current PA source, why a shorter emitter must be checked
logically before its state count is credited. This investigation did not
obtain such a new calibration or a smaller complete compilation from the
concise-iPRA presentations.

PA194P5 preserves the supplied logical calibration precisely
because it changes physical representation instead. The current
state count remains $134+67-7=194$. No global minimality, new
logical axiom system, or general runtime advantage is claimed.

\section{Reproduction and limitations}

From the extracted package directory run
\begin{verbatim}
python3 verify.py
python3 comparison/verify_novelty.py
python3 comparison/verify_census.py
\end{verbatim}
The second command replays the delivered frontier comparison; the third
reconstructs the complete delivered table census from its literal files.
For the optional bounded local searches:
\begin{verbatim}
cd research/local_reductions
node local_search.js
node local_search.js --extended
\end{verbatim}
The extended run is research code; it is not needed to establish
the final table or its physical transport.

The main article builds with
\begin{verbatim}
cd report
pdflatex -interaction=nonstopmode -halt-on-error report.tex
pdflatex -interaction=nonstopmode -halt-on-error report.tex
\end{verbatim}
The raw mathematical sources and literal tables are preserved.
The checksum manifest binds the delivered files. Archive and
table metadata record the exact pinned inputs.

The strict sub-194 objective remains unfulfilled. The positive
result is a distinct 194-state physical implementation with
a complete source-preservation argument and reproducible
certificates. Its consistency equivalence inherits the
parent packet's ordinary PA calibration. Its formal-proof
status has not been upgraded by the new numerical checks.
These qualifications delimit the actual result.

\appendix
\clearpage
\section{Complete PA194P5 transition table}
\label{app:table}
Each entry gives the write bit, direction and target. $H$ denotes
halt label 194. Start in state zero on the all-zero tape.
The table uses the parent's fixed dense numbering.
{\small\ttfamily
\begin{longtable}{@{}rll@{\hspace{1.4em}}rll@{\hspace{1.4em}}rll@{}}
\toprule State & Read 0 & Read 1 & State & Read 0 & Read 1 & State & Read 0 & Read 1\\\midrule\endhead
0 & 1R1 & 1R2 & 65 & 0R91 & 1R49 & 130 & 0R134 & 1R135\\
1 & 1R3 & 1R1 & 66 & 0R92 & 0L75 & 131 & 0R136 & 1R137\\
2 & 0R0 & 1R2 & 67 & 0R86 & 1R33 & 132 & 0R138 & 1R139\\
3 & 0R4 & 0R1 & 68 & 0R93 & 1R68 & 133 & 0R140 & 1R141\\
4 & 1R5 & 1L6 & 69 & 0L94 & 1L33 & 134 & 0R100 & 1R67\\
5 & 1L7 & 1L8 & 70 & 0L84 & 0R2 & 135 & 0R142 & 1R127\\
6 & 0L7 & 1L6 & 71 & 0R95 & 1R71 & 136 & 0R143 & 1R144\\
7 & 0L9 & 1L6 & 72 & 0R75 & 1R96 & 137 & 0R145 & 1R146\\
8 & 0L10 & 0L11 & 73 & 0R59 & 0R60 & 138 & 0R134 & 1R147\\
9 & 1R12 & 0L9 & 74 & 0R86 & 1R97 & 139 & 0R120 & 1R148\\
10 & 0L9 & 0L10 & 75 & 0R42 & 0L13 & 140 & 0R149 & 1L131\\
11 & 1L8 & 1L11 & 76 & 0R87 & 1R37 & 141 & 0R150 & 1R151\\
12 & 0R12 & 0R13 & 77 & 0R30 & 1R98 & 142 & 0R79 & 1R41\\
13 & 0L14 & 0R13 & 78 & 0R42 & 1R99 & 143 & 0R28 & 1R152\\
14 & 1L15 & 0L15 & 79 & 0R100 & 1R101 & 144 & 0L153 & 1L154\\
15 & 0L15 & 1L16 & 80 & 0R102 & 1L103 & 145 & 0L129 & 1R155\\
16 & 0L17 & 1L16 & 81 & 0R30 & 1R80 & 146 & 0R156 & 1R157\\
17 & 0L18 & 1L16 & 82 & 0R97 & 1R104 & 147 & 0R149 & 1R158\\
18 & 0L19 & 1L16 & 83 & 0R30 & 1R105 & 148 & 0R159 & 1R120\\
19 & 0L20 & 1L16 & 84 & 1R84 & 1L15 & 149 & 0R100 & 1R160\\
20 & 0L21 & 1L16 & 85 & 0R106 & 1R85 & 150 & 0R161 & 1R162\\
21 & 0L22 & 1L16 & 86 & 0R73 & 1R107 & 151 & 0R163 & 0L10\\
22 & 0L23 & 1L16 & 87 & 0L46 & 0R93 & 152 & 0R164 & 1R165\\
23 & 0R23 & 1R24 & 88 & 0R108 & 1R98 & 153 & 0R166 & 1R157\\
24 & 0R25 & 1R26 & 89 & 0R108 & 1R69 & 154 & 0R167 & 1R156\\
25 & 0R27 & 1R28 & 90 & 0R109 & 0L14 & 155 & 0R135 & 1R130\\
26 & 0R29 & 1R30 & 91 & 0R110 & 1L74 & 156 & 0R147 & 1R41\\
27 & 0R31 & 1RH & 92 & 0R111 & 1R112 & 157 & 0R134 & 1R168\\
28 & 0L32 & 1R33 & 93 & 0R113 & 1R93 & 158 & 0R41 & 1L169\\
29 & 0R34 & 1R35 & 94 & 0R50 & 1R103 & 159 & 0L170 & 1R171\\
30 & 0L36 & 1L37 & 95 & 0R114 & 1R95 & 160 & 0R91 & 1R64\\
31 & 0R38 & 1R39 & 96 & 0R30 & 1R49 & 161 & 0R172 & 1R170\\
32 & 0R40 & 1R41 & 97 & 0R48 & 1R61 & 162 & 0R173 & 1R174\\
33 & 0R42 & 1R43 & 98 & 0R87 & 1L115 & 163 & 0R175 & 1R176\\
34 & 0R44 & 1R45 & 99 & 0R104 & 0L14 & 164 & 0R177 & 1R178\\
35 & 0R46 & 1R47 & 100 & 0R36 & 1R116 & 165 & 0L179 & 1L180\\
36 & 0R48 & 1R49 & 101 & 0R115 & 1R78 & 166 & 0R149 & 1R140\\
37 & 0R50 & 1R51 & 102 & 0L46 & 0R113 & 167 & 0R165 & 1R181\\
38 & 0R52 & 1R42 & 103 & 0L45 & 1R117 & 168 & 0L130 & 1R120\\
39 & 0R53 & 1R54 & 104 & 0L94 & 1R118 & 169 & 0R182 & 1R135\\
40 & 0R55 & 1R56 & 105 & 0R119 & 0L14 & 170 & 0R183 & 1R145\\
41 & 0R57 & 1R58 & 106 & 0R71 & 1R106 & 171 & 0R158 & 1R130\\
42 & 0R59 & 0R85 & 107 & 0R120 & 1R30 & 172 & 0R139 & 1R165\\
43 & 0R61 & 0L14 & 108 & 0L46 & 0R106 & 173 & 0R158 & 1R184\\
44 & 0R62 & 1R63 & 109 & 0R70 & 0R68 & 174 & 0R185 & 1R91\\
45 & 0R55 & 1R64 & 110 & 0R59 & 0R68 & 175 & 0R31 & 1R53\\
46 & 0L65 & 1R66 & 111 & 0L65 & 1R121 & 176 & 0R186 & 1R175\\
47 & 0R36 & 1R67 & 112 & 0R122 & 1R123 & 177 & 0R134 & 1R142\\
48 & 0R59 & 0R114 & 113 & 0R14 & 1R113 & 178 & 0R32 & 1R187\\
49 & 0R69 & 0L14 & 114 & 0R68 & 1R114 & 179 & 0R144 & 1R188\\
50 & 0R70 & 0R71 & 115 & 0R124 & 1L67 & 180 & 0R189 & 1R190\\
51 & 0R72 & 0L49 & 116 & 0R125 & 1R126 & 181 & 0R41 & 1R191\\
52 & 0R73 & 1R54 & 117 & 0R120 & 0L14 & 182 & 0R168 & 1R41\\
53 & 0L74 & 0L15 & 118 & 0R94 & 1L98 & 183 & 0R192 & 1R120\\
54 & 0R75 & 0L10 & 119 & 0R127 & 1R127 & 184 & 0R153 & 1R146\\
55 & 0R76 & 1R77 & 120 & 0R70 & 0R114 & 185 & 0R165 & 1R180\\
56 & 0R78 & 1R79 & 121 & 0R128 & 1R129 & 186 & 0R96 & 1R115\\
57 & 0R80 & 1R81 & 122 & 0R32 & 1R36 & 187 & 0R193 & 1R159\\
58 & 0R82 & 1R83 & 123 & 0R31 & 1R79 & 188 & 0R169 & 1R154\\
59 & 0L84 & 0R1 & 124 & 0R59 & 0R95 & 189 & 0R158 & 1R192\\
60 & 0R85 & 1R60 & 125 & 0R102 & 1R117 & 190 & 0R179 & 1R182\\
61 & 0L76 & 1L86 & 126 & 0R82 & 1R81 & 191 & 0R182 & 1R41\\
62 & 0R87 & 0L7 & 127 & 0R120 & 1R120 & 192 & 0R149 & 1R168\\
63 & 0R88 & 1R89 & 128 & 0R130 & 1R131 & 193 & 0R79 & 1R144\\
64 & 0R42 & 1R90 & 129 & 0R132 & 1R133 &  &  & \\
\bottomrule
\end{longtable}}


\clearpage
\section{The 72-packet comparison ledger}
\label{app:ledger}
The placed-packet descriptions are attributed claims recorded
in the pinned placement ledger; the incoming descriptions are
read from those archives. Tables and graph identities were
audited separately. File-level provenance and duplicate mappings
are retained in the machine-readable census.
{\small\begin{longtable}{@{}>{\raggedright\arraybackslash}p{0.36\linewidth}>{\raggedright\arraybackslash}p{0.59\linewidth}@{}}
\toprule Packet & Reported contribution\\\midrule\endhead
\multicolumn{2}{@{}l}{\textbf{Placed packets (61)}}\\\midrule
\nolinkurl{pa_unary_power_cover} & Unary power covering at PA strength (report of 27 October 2026). The U208 packets carry it as a baseline.\\[4pt]
\nolinkurl{U208_ConPA_Conditional_Context_2026-10-29} & U208 by a conditional tape-context reduction. States a 208-state candidate \nolinkurl{machine/u208.tm}.\\[4pt]
\nolinkurl{U208_ConPA_Entry_Context_2026-10-29} & U208 by eliminating a state using entry context. States a 208-state machine \nolinkurl{machine/u208.tm}.\\[4pt]
\nolinkurl{U208_ConPA_History_2026-10-29} & U208 by a finite-history decoder fold. States a 208-state candidate \nolinkurl{machine/u208.tm}.\\[4pt]
\nolinkurl{U208_ConPA_2026-10-29} & U208 by a three-cell fold. States a 208-state candidate \nolinkurl{machine/u208.tm}.\\[4pt]
\nolinkurl{u208_research_2026-10-29} & U208 research release. States two 208-state machines, \nolinkurl{machine/u208.tm} and \nolinkurl{machine/variant_plain/u208.tm}.\\[4pt]
\nolinkurl{Below_U208_Research_2026-10-29} & Research toward fewer than 208 states. States that no machine with 207 states or fewer was obtained.\\[4pt]
\nolinkurl{u208_comparison_research_2026-10-29} & Comparison of the 208-state proposals and a search below 208. States that the target was not achieved.\\[4pt]
\nolinkurl{ConPA_below208_research_2026-10-29} & Comparative audit and reduction search below 208. States that it does not establish a 207-state machine.\\[4pt]
\nolinkurl{pa_consistency_research_2026-10-29} & Audit and new source reductions. States that no validated machine below 208 states was obtained.\\[4pt]
\nolinkurl{pa208_research_2026-10-29} & Comparison and research below 208 states. States that it does not establish a valid machine below 208.\\[4pt]
\nolinkurl{pa_tm_research_2026-11-13} & Comparison and reduction research. States a new 196-state variant with shorter source paths and nothing below 196.\\[4pt]
\nolinkurl{pa_consistency_research_2026-11-13} & Comparison and diagonal-vacuity research. States a shorter axiom program with a 199-state implementation and nothing below 196.\\[4pt]
\nolinkurl{conpa_tm_comparison_and_reduction_search} & Comparison and reduction search. States two experimental 196-state command variants and nothing below 196.\\[4pt]
\nolinkurl{pa196_research_2026-11-13} & PA196 comparison, source factorization and proof obligations. States a source tree of 157 internal nodes and nothing below 196.\\[4pt]
\nolinkurl{pa_consistency_research} & Research and verification packet. States that no machine below 196 states has been validated.\\[4pt]
\nolinkurl{pa_tm_comparison_and_reductions} & Comparison and new rewrites, a research snapshot.\\[4pt]
\nolinkurl{pa_tm_extended_research_2026-11-13} & Comparison and refined reduction search. States that the smallest accepted count is 196.\\[4pt]
\nolinkurl{conpa_tm_comparison_and_search} & Comparison and a search below 196. States that none was obtained; files named \nolinkurl{195} are rejected or diagnostic candidates.\\[4pt]
\nolinkurl{pa_consistency_reduction} & Comparison and reduction experiments. States two new 196-state variants and a shorter source program with a 203-state implementation.\\[4pt]
\nolinkurl{pa_consistency_comparison_2026_11_13} & Comparison and new reduction experiments. States that nothing below 196 was obtained.\\[4pt]
\nolinkurl{PA_consistency_TM_research_2026-11-13} & Comparison and new source reductions. States that no machine below 196 is claimed.\\[4pt]
\nolinkurl{pa_consistency_research_2026-11-13_2} & Comparison and search below 196. States two new layouts with 196 and 198 states and a command-level reduction.\\[4pt]
\nolinkurl{conpa_tm_study_2026-11-13} & Comparison and reduction study. States that nothing below 196 was obtained.\\[4pt]
\nolinkurl{pa_below196_research} & Comparison and reduction experiments. States that nothing below 196 was obtained.\\[4pt]
\nolinkurl{pa_consistency_comparison_and_new196_2026-11-13} & Comparison and reduction research. States that it contains no accepted machine below 196.\\[4pt]
\nolinkurl{pa_tm_research_2026-11-13_2} & Comparison and reduction investigation. States same-size work, failed candidates and nothing below 196.\\[4pt]
\nolinkurl{pa_consistency_cross_archive} & Cross-archive constructions and reduction search. States that it does not certify a 195-state machine.\\[4pt]
\nolinkurl{pa195_conpa_2026-11-13} & PA195. States a 195-state machine \nolinkurl{pa195/pa195.tm} and an article with a PA proof of the equivalence.\\[4pt]
\nolinkurl{pa_consistency_reduction_audit_2026-11-14} & Comparison and reduction audit. States nothing below 196 and no new theorem.\\[4pt]
\nolinkurl{pa_consistency_research_2026-11-14} & Comparison and NEG18. States an 18-instruction encoded negation with a compiled machine of 197 states.\\[4pt]
\nolinkurl{PA195_ConPA_research_2026-11-14} & A 195-state machine. States the primary \nolinkurl{machine_195S.tm} and a second table \nolinkurl{machine_195C.tm}.\\[4pt]
\nolinkurl{pa195_followup_research} & Comparison and searches below 195. States a graph-distinct 195-state machine M195S8 and nothing below 195.\\[4pt]
\nolinkurl{pa194_conpa_2026-11-13} & A 194-state machine. States the table \nolinkurl{construction/tables/pa194.tm} and a proof of the equivalence in PA.\\[4pt]
\nolinkurl{pa194_conpa_2026-11-13_2} & PA194. States a 194-state machine \nolinkurl{pa194/pa194.tm} and an article with a PA proof of the equivalence.\\[4pt]
\nolinkurl{pa195_diagonal_combined_2026-11-13} & PA195DC. States a different 195-state machine.\\[4pt]
\nolinkurl{pa194C_conpa_2026-11-13_2} & PA194C. States a materially different 194-state machine, an article with a PA proof of the equivalence, and nothing below 194.\\[4pt]
\nolinkurl{pa194_new_designs_2026-11-13} & Two new 194-state machines, PA194C and PA194CS. States that nothing below 194 was obtained.\\[4pt]
\nolinkurl{pa194c_conpa} & PA194C, a cached-head 194-state machine. States that no machine below 194 was established.\\[4pt]
\nolinkurl{pa194c_conpa_2026-11-13_1} & PA194C. States a different 194-state machine and an article with a PA proof of the equivalence.\\[4pt]
\nolinkurl{pa194c_conpa_2026-11-13} & PA194C, a cached-head 194-state machine. States an article with a PA proof of the equivalence and nothing below 194.\\[4pt]
\nolinkurl{pa194_radix_family_2026-11-13} & New 194-state machines calibrated to Con(PA), a radix family.\\[4pt]
\nolinkurl{pa194j_pa194d_conpa_2026-11-13} & PA194J and PA194D. States two new 194-state machines.\\[4pt]
\nolinkurl{pa194q_conpa_2026-11-13} & PA194Q and PA194CSZ. States two new 194-state machines and that no machine below 194 was established.\\[4pt]
\nolinkurl{PA194R_research_2026-11-14} & PA194R. States a different 194-state machine with a new decoder layout and that no machine below 194 was verified.\\[4pt]
\nolinkurl{pa194N_conpa_2026-11-13} & PA194N. States a new 194-state machine, an article with a PA proof of the equivalence, and nothing below 194.\\[4pt]
\nolinkurl{pa194zr_conpa_2026-11-13} & PA194ZR. States a new 194-state machine \nolinkurl{machine/pa194zr.tm}, a PA proof of the equivalence, and nothing below 194.\\[4pt]
\nolinkurl{pa194hs_conpa_2026-11-13} & PA194HS and PA194H. States two 194-state machines and that no machine below 194 was established.\\[4pt]
\nolinkurl{pa194n5_n6_conpa_2026-11-13} & PA194N5 and PA194N6. States two distinct 194-state machines and claims no machine below 194.\\[4pt]
\nolinkurl{pa194ns_conpa_2026-11-13} & PA194NS. States a different 194-state machine and that no machine below 194 was established.\\[4pt]
\nolinkurl{pa194n6_research_2026-11-13} & PA194N6. States a different 194-state machine and that no machine below 194 was verified.\\[4pt]
\nolinkurl{PA194NS_research} & PA194NS and PA194SUMS. States two new 194-state constructions.\\[4pt]
\nolinkurl{pa194haz_research_2026-11-14} & PA194HA and PA194HAZ. States two 194-state alternatives and that no machine below 194 was obtained.\\[4pt]
\nolinkurl{pa194qn_research_2026-11-13} & PA194QN. States a distinct 194-state machine and that no verified machine below 194 was found.\\[4pt]
\nolinkurl{pa194h6_conpa_2026-11-13} & PA194H6. States a new 194-state machine, an article with a PA proof of the equivalence, and nothing below 194.\\[4pt]
\nolinkurl{pa194dhs_conpa_2026-11-13} & PA194DHS. States a different 194-state machine and that no machine below 194 was verified.\\[4pt]
\nolinkurl{pa194nq_research_2026-11-13} & PA194NQ. States a new 194-state machine and claims no machine below 194.\\[4pt]
\nolinkurl{pa194_qns_research_2026-11-13} & PA194QNS. States a new 194-state machine and that no verified construction below 194 was found.\\[4pt]
\nolinkurl{pa194_qradix_research_2026-11-13} & PA194Q5 and PA194Q6. States two different 194-state machines and that nothing below 194 was obtained.\\[4pt]
\nolinkurl{pa194qh6_conpa_2026-11-13} & PA194QH6. States a new 194-state composition and that no machine below 194 was obtained.\\[4pt]
\nolinkurl{pa194qns_research_2026-11-13} & PA194QNS. States a distinct 194-state machine and that no machine below 194 was established.\\[4pt]
\midrule\multicolumn{2}{@{}l}{\textbf{Current incoming archives (11)}}\\\midrule
\nolinkurl{pa194_ax11_research_2026-11-14} & PA194AX11 changes the AX command to eleven stack entries and four encoded operands by falling through into the full MP reader, while removing the MP dummy push; also supplies QH5. Both 194. Ordinary PA proofs and physical certificates; no sub-194 result.\\[4pt]
\nolinkurl{pa194_qh5_research_2026-11-14} & PA194QH5 composes coefficient five, exchanged quantifier premises, no-push MP and cached-only Shift; duplicates the QH5 table in AX11. Includes failed runtime193 startup repairs, cyclic mergers and an eager-Shift counterexample.\\[4pt]
\nolinkurl{pa194h5_research_2026-11-14} & PA194H5 specializes H to coefficient five; new composition, not a new search architecture. Also PA195EvenHX, a 195-state scheduler enumerating even candidates with a completeness proof. No sub-194 result.\\[4pt]
\nolinkurl{pa194qh_conpa_2026-11-13} & PA194QH combines exchanged quantifier premises, visible no-push MP and cached-only Shift; ordinary PA transfer plus symbolic physical certificates. No sub-194 result.\\[4pt]
\nolinkurl{pa194qhs_conpa_2026-11-14} & PA194QHS adds symmetric cache-only Shift to QH, eliminating a restoration loop with unbounded source-instruction savings. Ordinary PA proof and physical certificates; no sub-194 result.\\[4pt]
\nolinkurl{pa194qn6_conpa_2026-11-14} & PA194QN6 combines exchanged quantifier premises, coefficient six and no-push MP while retaining sum-based Shift. Same principal composition appears in the QN5/QN6 packet. No sub-194 result.\\[4pt]
\nolinkurl{pa194qn_2026-11-14} & PA194QN5 and companion PA194QN6 combine exchanged quantifier premises, no-push MP, sum-based Shift and coefficient five or six. Constituent ideas inherited; exact constructions and bindings supplied. No sub-194 result.\\[4pt]
\nolinkurl{pa_consistency_reduction_study_2026-11-13} & Comparative audit of 61 packets / 49 frontier classes; exact rational obstruction for a frozen helper-selector family, independently replayed coefficient-three refutation, and source-only reassociation checks. No accepted sub-194 machine.\\[4pt]
\nolinkurl{pa_consistency_research_2026-11-14} & Comparison and theorem excluding every nonnegative constant equality offset in the specified otherwise unchanged coefficient-three calculus; source alternatives and bounded physical searches. No accepted sub-194 machine.\\[4pt]
\nolinkurl{pa_sub194_research} & MP196 halts immediately when a successful MP consequent is one; a 196-state prototype with revised command invariant. Comparative audit and failed bounded reduction searches; no sub-194 result.\\[4pt]
\nolinkurl{pa_sub194_research_2026-11-13} & Runtime193 folds native 301 into 88 and changes two zero transitions. All 4,780 source contracts hold; 1158 elementary unchanged startups are rejected. A 193-state runtime core, not a calibrated blank-start consistency machine.\\[4pt]
\bottomrule
\end{longtable}
}

\clearpage
\begin{thebibliography}{99}
\bibitem{status} ProveIt, \emph{PA consistency machines}, pinned README.
\url{https://github.com/VladimirReshetnikov/ProveIt/blob/4be5c6e3a10292a7466f5d3687e3b39615ff049b/Computability/BusyBeaver/Research/PA/README.md}.
\bibitem{primary} \emph{A 194-state Turing machine calibrated to the consistency of PA},
primary delivered packet, 13 November 2026.
\url{https://github.com/VladimirReshetnikov/ProveIt/tree/4be5c6e3a10292a7466f5d3687e3b39615ff049b/Computability/BusyBeaver/Research/PA/docs/pa194_conpa_2026-11-13}.
\bibitem{calibration} Primary packet, \file{report/calibration.tex}.
A byte-preserved copy is \file{reference/primary_calibration.tex} in this package.
\bibitem{catalog} ProveIt, \emph{Machines in the delivered packets},
\file{PA/docs/MACHINES.md} and \file{MACHINES.tsv}, pinned revision.
\url{https://github.com/VladimirReshetnikov/ProveIt/blob/4be5c6e3a10292a7466f5d3687e3b39615ff049b/Computability/BusyBeaver/Research/PA/docs/MACHINES.md}.
\bibitem{placement} ProveIt, \emph{Delivered research packets on Con(PA) machines},
\file{PA/docs/README.md}, pinned revision; retained in the comparison directory.
\bibitem{rejected193} NQ research packet,
\file{search/rejected193/receipt.json}, pinned revision.
\url{https://github.com/VladimirReshetnikov/ProveIt/blob/4be5c6e3a10292a7466f5d3687e3b39615ff049b/Computability/BusyBeaver/Research/PA/docs/pa194nq_research_2026-11-13/search/rejected193/receipt.json}.
\bibitem{runtime193} Incoming \file{pa_sub194_research_2026-11-13.zip},
\file{runtime193/check_runtime88301.py} and its report.
\url{https://github.com/VladimirReshetnikov/ProveIt/blob/4be5c6e3a10292a7466f5d3687e3b39615ff049b/docs/incoming/pa_sub194_research_2026-11-13.zip}.
\bibitem{zfc} ProveIt, \emph{Smaller machines for Con(ZFC)}, pinned README.
\url{https://github.com/VladimirReshetnikov/ProveIt/blob/4be5c6e3a10292a7466f5d3687e3b39615ff049b/Computability/BusyBeaver/Research/ZFC/README.md}.
\bibitem{concise} ProveIt, \emph{Concise presentations of intuitionistic primitive recursive arithmetic}, pinned README.
\url{https://github.com/VladimirReshetnikov/ProveIt/blob/4be5c6e3a10292a7466f5d3687e3b39615ff049b/Logic/PrimitiveRecursiveArithmetic/ConciseIPRA/README.md}.
\bibitem{boxlemmas} ProveIt / James Chang, \file{M6.Native.boxLemmas},
append-one arithmetic-box proof, pinned source.
\url{https://github.com/VladimirReshetnikov/ProveIt/blob/4be5c6e3a10292a7466f5d3687e3b39615ff049b/Computability/BusyBeaver/Research/PA/jamesyc-upstream/pa196/M6/Native.lean}.
\end{thebibliography}
\end{document}
